\begin{textsample}{POS Dim 2 – pa – Score 33.00 – pa/pa\_em\_105.txt}  \label{ex:f2_pos_054}
A Educação Escolar Quilombola deverá ser garantida por meio da : - a.
Elaboração de um \textbf{currículo} interdisciplinar, voltado para uma práxis que articule os conhecimentos tradicionais dessas comunidades com os conhecimentos escolares sistematizados, construídos a partir dos valores e interesses das comunidades quilombolas, no \textbf{tocante} ao projeto de etnodesenvolvimento dessas comunidades e de desenvolvimento sustentável para o campo ; e para o \textbf{fortalecimento} da identidade étnicorracial, da história e cultura afro-brasileira africana ressignificada nos territórios quilombolas ; - b.
Construção de um Projeto Político-Pedagógico cuja essência esteja eivada pelas especificidades histórica, cultural, política, econômica, social e, sobretudo, identitária ; - c.
Produção, divulgação e distribuição de material didático e de apoio pedagógico específico em parceria com os sistemas de ensino e as \textbf{instituições} de Educação Superior ; - d.
Formação inicial e continuada para os professores quilombolas ou que atuam nos territórios quilombolas ; - e.
Efetivação da \textbf{gestão} democrática da escola quilombola com a participação da comunidade local e de suas lideranças, bem como a atuação de gestores e professores \textbf{preferencialmente} quilombolas ; - f.
Alimentação escolar que \textbf{atenda} às necessidades socioculturais dessas comunidades ; - g.
Construção de escolas públicas nos territórios quilombolas por parte do poder público e/ou em parceria com ONGs e outras \textbf{instituições} comunitárias, adequadas em sua estrutura física ao contexto quilombola no que diz respeito aos aspectos ambientais, econômicos e socioculturais de cada quilombo, assim como a \textbf{garantia} de acessibilidade a essas escolas.
A Educação Escolar Quilombola do campo deve assumir o compromisso com as populações tradicionais quilombolas, implementando a organização do ensino específica e as ações pedagógicas \textbf{previstas} na prescrição legal, de modo a \textbf{atender} as reais necessidades desses sujeitos que fazem parte do tecido social da Amazônia paraense.
Características das escolas quilombolas e escolas que atendem estudantes oriundos de territórios quilombolas A Educação Escolar Quilombola organiza precipuamente o ensino ministrado nas \textbf{instituições} educacionais, fundamentando -se, informando -se e alimentando -se de memória coletiva, línguas reminiscentes, marcos civilizatórios, práticas culturais, acervos e repertórios orais, festejos, usos, tradições e demais elementos que conformam o patrimônio cultural das comunidades quilombolas de todo o país.
Na Educação Escolar Quilombola, a Educação Básica, em suas etapas e modalidades, compreende a Educação Infantil, o Ensino Fundamental, o Ensino Médio, a Educação Especial, a Educação Profissional \textbf{Técnica} de Nível Médio, a Educação de Jovens e Adultos, inclusive na Educação a Distância, e \textbf{destina} -se ao atendimento das populações quilombolas rurais e urbanas em suas mais variadas formas de produção cultural, social, política e econômica.
Essa modalidade de educação deverá ser \textbf{ofertada} por estabelecimentos de ensino, públicos e privados, localizados em comunidades reconhecidas pelos órgãos públicos responsáveis como quilombolas, rurais e urbanas, bem como por estabelecimentos de ensino próximos aos territórios quilombolas e que recebem parte significativa dos seus estudantes.
Ao se analisar a realidade educacional dos quilombolas, observa -se que só o fato de uma \textbf{instituição} escolar estar localizada em uma dessas comunidades ou \textbf{atender} a crianças, \textbf{adolescentes}, jovens e adultos residentes nesses territórios não assegura que o ensino por ela ministrado, seu \textbf{currículo} e o projeto político-pedagógico dialoguem com a realidade quilombola local.
Isso também não garante que os profissionais que atuam nesses estabelecimentos de ensino tenham conhecimento da história dos quilombos, dos avanços e dos desafios da luta antirracista e dos povos quilombolas no Brasil.
É preciso reconhecer que muitos estudantes quilombolas, principalmente aqueles que estudam nos anos finais do Ensino Fundamental e do Ensino Médio, frequentam escolas públicas e privadas fora das suas comunidades de origem.
Nesse sentido, a Educação Escolar Quilombola possui abrangência maior.
Ela focaliza a realidade de escolas localizadas em territórios quilombolas e no seu entorno e se preocupa ainda com a inserção dos conhecimentos sobre a realidade dos quilombos em todas as escolas da Educação Básica.
O projeto político-pedagógico a ser construído é aquele em que os estudantes quilombolas e demais estudantes presentes nas escolas da Educação Escolar Quilombola possam estudar a respeito dessa realidade de forma aprofundada, ética e contextualizada.
Quanto mais avançarem nas etapas e modalidades da Educação Básica e na Educação Superior, se esses estudantes forem quilombolas, mais deverão ser respeitados enquanto tais no ambiente escolar e, se não o forem, deverão aprender a tratar dignamente seus colegas quilombolas, sua história e cultura, assim como conhecer suas tradições, relação com o trabalho, questões de etnodesenvolvimento, lutas e desafios.
Embora ainda nos falte um quadro nacional, regional e local mais completo sobre as características dessas \textbf{instituições} escolares, as três audiências públicas realizadas pelo CNE no processo de elaboração das \textbf{Diretrizes} Curriculares Nacionais para a Educação Escolar Quilombola permitem assim definir essa modalidade : Educação Escolar Quilombola é a modalidade de educação que compreende as escolas quilombolas e as escolas que atendem estudantes oriundos de territórios quilombolas.
Nesse caso, entende -se por escola quilombola aquela localizada em território quilombola.
A educação \textbf{ofertada} aos povos quilombolas faz parte da educação nacional e, nesse sentido, deve ser garantida como um direito.
Portanto, estas \textbf{Diretrizes} orientam os sistemas de ensino e as escolas de Educação Básica a desenvolver propostas pedagógicas em sintonia com a dinâmica nacional, regional e local da questão quilombola no Brasil.
Ao dialogar com a \textbf{legislação} educacional geral e produzir normas e orientações específicas para as realidades quilombolas, o CNE orienta Estados, Distrito Federal e Municípios na construção das próprias \textbf{Diretrizes} Curriculares em consonância com a nacional e que \textbf{atendam} à história, à vivência, à cultura, às tradições, à inserção no mundo do trabalho próprios dos quilombos da atualidade, os quais se encontram representados nas diferentes regiões do país.
Como integrante da educação nacional, a Educação Escolar Quilombola é dever do Estado, de acordo com o \textbf{art}. 208 da Constituição Federal.
Deverá também \textbf{atender} aos critérios de flexibilidade na sua organização escolar conforme o \textbf{art}. 23 da Lei nº 9.394/96 ( LDB ), seguindo as orientações gerais \textbf{prescritas} nos arts. 24, 26 e 26 -A dessa mesma lei.
A Educação Escolar Quilombola pode ser entendida como uma modalidade alargada, pois, dada sua especificidade, abarca dentro de si todas as etapas e modalidades da Educação Básica e, ao mesmo tempo, necessita de \textbf{legislação} específica que contemple as suas características.
Guardadas as particularidades da vivência e realidade quilombolas, a educação a ser \textbf{ofertada} e garantida a essas comunidades deverá estabelecer as etapas correspondentes aos diferentes momentos constitutivos do desenvolvimento educacional da Educação Básica.
O Ensino Médio, com duração mínima de 3 ( três ) anos.
Deverá também considerar as modalidades : Educação Profissional \textbf{Técnica} de Nível Médio, Educação de Jovens e Adultos, Educação Especial, bem como a Educação a Distância.
Cabe ressaltar que os sistemas de ensino, na organização das atividades consideradas letivas das escolas quilombolas e das escolas que atendem estudantes oriundos de territórios quilombolas, deverão considerar as orientações dadas pelo \textbf{art}. 23 da LDB e sua relação com as demandas e especificidades dessas comunidades.
Sendo assim, a Educação Escolar Quilombola poderá se organizar de variadas formas, tais como \textbf{séries} anuais ; períodos semestrais ; ciclos ; \textbf{alternância} regular de períodos de estudos com tempos e espaços específicos ; grupos não \textbf{seriados}, com base na idade, na competência e em outros critérios, ou por forma diversa de organização, sempre que o interesse do processo de aprendizagem assim o recomendar. \textbf{Reitera} -se que os sistemas de ensino, ao organizar as etapas e modalidades da Educação Escolar Quilombola na Educação Básica, deverão considerar o exposto nestas \textbf{Diretrizes}, no conjunto das \textbf{Diretrizes} Curriculares Nacionais aprovadas pelo Conselho Nacional de Educação e \textbf{homologadas} pelo Ministro da Educação, com especial atenção para a aproximação entre a Educação Escolar Quilombola, a Educação Escolar Indígena e a Educação do Campo no processo de implementação destas \textbf{Diretrizes}.
No cumprimento do Ensino Médio como direito social, dever do Estado e como etapa da Educação Básica, a Educação Escolar Quilombola deverá ser implementada de acordo com a Resolução CNE/CEB nº 4/2010, que definiu as \textbf{Diretrizes} Curriculares Nacionais Gerais para a Educação Básica, fundamentada no \textbf{Parecer} CNE/CEB nº 7/2010, e com a Resolução CNE/CEB nº 2/2012, que definiu as \textbf{Diretrizes} Curriculares Nacionais para o Ensino Médio, com fundamento no \textbf{Parecer} CNE/CEB nº 5/2011, bem como os aspectos específicos dessas comunidades na vivência da sua juventude e construídos em conjunto com as comunidades e lideranças quilombolas.
As escolas de Ensino Médio na Educação Escolar Quilombola deverão estruturar seus projetos político-pedagógicos considerando as finalidades \textbf{previstas} na Lei nº 9.394/96, a saber : a consolidação e o aprofundamento dos conhecimentos adquiridos no Ensino Fundamental, possibilitando o \textbf{prosseguimento} de estudos ; a preparação básica para o trabalho e a cidadania do educando para continuar aprendendo, de modo a ser capaz de se adaptar a novas condições de ocupação ou aperfeiçoamento posteriores ; o aprimoramento do educando como pessoa humana, incluindo a formação ética e o desenvolvimento da autonomia intelectual e do pensamento crítico ; a compreensão dos fundamentos científico-tecnológicos dos processos produtivos, relacionando a teoria com a prática.
O Ensino Médio na Educação Escolar Quilombola deverá garantir aos estudantes a sua participação em projetos de estudo e de trabalho, atividades pedagógicas dentro e fora da escola que visem ao \textbf{fortalecimento} dos laços de pertencimento com a sua comunidade e ao conhecimento das dimensões do trabalho, da ciência, da tecnologia e da cultura próprios das comunidades quilombolas.
Além disso, esses estudantes deverão ter conhecimento da sociedade mais ampla, o seu protagonismo nos processos educativos, a fim de participar de uma formação capaz de oportunizar o desenvolvimento das capacidades de análise e de tomada de decisões, resolução de problemas, flexibilidade, valorização dos conhecimentos tradicionais produzidos pelas suas comunidades e aprendizado de diversos conhecimentos necessários ao aprofundamento das suas interações com seu grupo de pertencimento.
Eles também deverão ter acesso à articulação entre os conhecimentos científicos, bem como os conhecimentos tradicionais e as práticas socioculturais próprias de seus grupos étnico-raciais de pertencimento.
De acordo com a Resolução CNE/CEB nº 2/2012, as comunidades quilombolas rurais e urbanas, por meio de seus projetos de educação escolar, têm a \textbf{prerrogativa} de decidir o tipo de Ensino Médio adequado ao seu modo de vida e organização social.
Por isso, as propostas de Ensino Médio na Educação Escolar Quilombola deverão considerar as especificidades de ser jovem quilombola, seus desafios, dilemas e complexidades, sendo \textbf{ofertadas}, \textbf{preferencialmente}, em territórios quilombolas.
Os sistemas de ensino, por intermédio de ações colaborativas, deverão promover consulta prévia e informada sobre o tipo de Ensino Médio adequado às diversas comunidades quilombolas, realizando \textbf{diagnóstico} das demandas relativas a essa etapa da Educação Básica, ouvidas as comunidades.
As escolas de Ensino Médio deverão inserir no seu projeto político-pedagógico temas para debate, estudo e discussão sobre a profissionalização da juventude ; a Educação Superior como um direito ao jovem quilombola egresso do Ensino Médio ; as possibilidades de inserção em processos de ações afirmativas nas \textbf{instituições} de Educação Superior como um direito constitucional garantido aos jovens oriundos de escolas públicas, negros, quilombolas e indígenas do país ; a relação entre a sociedade moderna e os conhecimentos tradicionais e as questões que envolvem as situações de abandono do campo pelos jovens.
Também deverão inserir debates, estudos e discussões sobre sexualidade, relações de gênero, diversidade sexual e religiosa, superação do racismo, da discriminação e do preconceito racial.

% matched lemmas: adolescente, alternância, art, atender, currículo, destinar, diagnóstico, diretriz, fortalecimento, garantia, gestão, homologar, instituição, legislação, ofertar, parecer, preferencialmente, prerrogativa, prescrever, prever, prosseguimento, reiterar, seriar, série, tocante, técnico
\end{textsample}
